\documentclass[UTF8,12pt,a4paper]{ctexart}

\usepackage{amsmath,amssymb}
\usepackage{booktabs}
\usepackage{geometry}
\usepackage{graphicx}
\usepackage[hidelinks]{hyperref}
\geometry{left=2.5cm,right=2.5cm,top=2.5cm,bottom=2.5cm}
\newcommand{\keywords}[1]{\par\noindent\textbf{关键词：}#1}

% ============================================================
% 国赛官方格式要点（2026 年修订稿，以当届官网为准）：
%   - 纸质版：第一页承诺书、第二页编号专用页、第三页摘要专用页
%     （摘要含标题和关键词、无需翻译成英文，不超一页；页码从摘要页
%     "1"开始、位于页脚中部）。
%   - 正文从第四页开始：不要目录，不超过 30 页；正文后附录页数不限。
%   - 附录必须含支撑材料文件列表 + 全部完整可运行源程序；
%     确实没有用到程序须注明"本论文没有用到程序"。
%   - 电子版：第一页必须是摘要专用页（不要承诺书/编号页），
%     PDF/Word 单文件 ≤20MB、不压缩；全文不得出现身份/学校/赛区信息。
%   - 当届官方模板发布后必须优先使用并核验；本文件只是构建基线，
%     不替代官方文件。
% ============================================================

\title{LATEX_TEMPLATE_TITLE}
\author{}
\date{}

\begin{document}

% 纸质版需在摘要页之前插入官方承诺书页与编号专用页（本基线不含）。
% 电子版第一页即本页（摘要专用页），页码从 1 开始。
\begin{abstract}
LATEX_TEMPLATE_ABSTRACT
\keywords{LATEX_TEMPLATE_KEYWORDS}
\end{abstract}

\section{问题重述}
LATEX_TEMPLATE_BODY

\section{问题分析}
LATEX_TEMPLATE_BODY

\section{模型假设与符号说明}
LATEX_TEMPLATE_BODY

\section{数据预处理}
LATEX_TEMPLATE_BODY

% 逐问建模与求解：按题目子问题数量展开，每问一节；子问题只有两个时删除多余小节。
% 每问包含：建模思路 → 数学表达 → 求解方法 → 结果 → 图表 + 紧随其后的分析。
\section{问题一建模与求解}
LATEX_TEMPLATE_BODY

\section{问题二建模与求解}
LATEX_TEMPLATE_BODY

\section{问题三建模与求解}
LATEX_TEMPLATE_BODY

\section{模型的分析与检验}
LATEX_TEMPLATE_BODY

\section{模型的评价与推广}
LATEX_TEMPLATE_BODY

\bibliographystyle{plain}
\bibliography{references}

\appendix
\section{附录}
% 附录：支撑材料文件列表 + 全部完整可运行源程序。
% 没有用到程序时在此注明"本论文没有用到程序"。
LATEX_TEMPLATE_BODY

\end{document}
